\chapter{Antecedentes}\label{chapter:antecedentes}

\section{Marco Te\'orico}

\subsection{Gesti\'on de establecimientos bovinos de cr\'ia}

La gesti\'on de un establecimiento bovino de cr\'ia requiere administrar informaci\'on productiva, reproductiva, sanitaria y econ\'omica de manera integrada. En este tipo de explotaciones, cada animal constituye una unidad de seguimiento relevante, dado que su historial de peso, edad, raza, vacunaciones, tratamientos, sangrados, tactos y estado reproductivo incide en decisiones de manejo y comercializaci\'on.

En la pr\'actica, parte de esta informaci\'on puede registrarse en cuadernos, planillas de c\'alculo o sistemas no integrados. Si bien estos mecanismos pueden resultar suficientes para establecimientos peque\~nos o con baja complejidad operativa, presentan limitaciones cuando aumenta la cantidad de animales, usuarios intervinientes o actividades sanitarias. MiGanado toma como referencia un escenario operativo en el cual un establecimiento puede administrar rodeos de tama\~no medio, con m\'ultiples actores involucrados y necesidad de consulta frecuente, lo que vuelve necesario contar con mecanismos de registro, actualizaci\'on, trazabilidad interna y control m\'as eficientes.

La digitalizaci\'on de estos procesos no implica solamente reemplazar el soporte papel por una pantalla. Tambi\'en exige definir entidades, reglas de negocio, permisos de acceso, criterios de calidad de datos y mecanismos de auditor\'ia. En consecuencia, el problema abordado por el proyecto combina aspectos del dominio ganadero con decisiones propias de ingenier\'ia de software: modelado de datos, dise\~no de interfaces, seguridad, validaci\'on funcional y mantenimiento evolutivo.

\subsection{Identificaci\'on animal y caravana}

La caravana es el identificador asociado al animal y colocado en su oreja, que permite reconocerlo de manera individual dentro del rodeo. En el contexto de MiGanado, el n\'umero de caravana funciona como clave operativa para consultar y registrar eventos sanitarios, reproductivos, productivos y econ\'omicos. El uso de este identificador facilita la trazabilidad interna y reduce la ambig\"uedad al momento de cargar informaci\'on.

En una versi\'on futura, el sistema podr\'ia integrarse con dispositivos de lectura externa, como bastones Bluetooth o mecanismos asociados a caravanas electr\'onicas. Sin embargo, de acuerdo con el alcance definido para el prototipo, dichos dispositivos no forman parte del desarrollo del proyecto y se consideran fuentes externas de datos.

\subsection{Ganader\'ia de precisi\'on}

La ganader\'ia de precisi\'on, o \emph{Precision Livestock Farming}, comprende el uso de tecnolog\'ias orientadas al monitoreo y an\'alisis de informaci\'on productiva, ambiental y sanitaria en sistemas ganaderos. Su prop\'osito es asistir a productores y profesionales mediante datos m\'as oportunos, automatizaci\'on de procesos y detecci\'on temprana de situaciones que podr\'ian afectar la productividad o el bienestar animal.

Schillings, Bennett y Rose analizan el potencial de estas tecnolog\'ias para abordar aspectos vinculados con el bienestar animal, destacando la importancia del monitoreo y de la informaci\'on sistem\'atica en establecimientos productivos \parencite{SchillingsBennettRose2021}. En l\'inea con este enfoque, MiGanado busca transformar registros dispersos en informaci\'on estructurada y accionable.

Desde la perspectiva de la agricultura digital, la incorporaci\'on de tecnolog\'ias de informaci\'on permite recolectar, almacenar, analizar y compartir datos a lo largo de los procesos productivos \parencite{FAO2019DigitalTechnologies}. Para este proyecto, dicho enfoque resulta relevante porque el valor de la plataforma depende de la continuidad del registro y de la posibilidad de convertir eventos cotidianos del establecimiento en indicadores comparables.

\subsection{Inteligencia artificial aplicada al sector ganadero}

La inteligencia artificial aplicada al sector ganadero permite analizar datos hist\'oricos, identificar patrones y asistir en la generaci\'on de estimaciones. En el proyecto MiGanado, estas capacidades se orientar\'an inicialmente a predicciones de peso, probabilidad de pre\~nez, alertas de salud, detecci\'on de anomal\'ias y proyecciones econ\'omicas.

No obstante, para que los resultados sean confiables, ser\'a necesario documentar la procedencia de los datos utilizados, los criterios de entrenamiento, las limitaciones de los modelos y las condiciones bajo las cuales las recomendaciones del sistema deben ser interpretadas. Esta precauci\'on es especialmente importante en un dominio donde una alerta autom\'atica puede orientar una decisi\'on productiva o sanitaria, pero no debe reemplazar el criterio de profesionales ni la observaci\'on directa del rodeo.

\subsection{Trazabilidad, calidad de datos y toma de decisiones}

La trazabilidad interna requiere que los eventos registrados puedan asociarse de manera consistente con animales, lotes, establecimientos, usuarios y fechas. Esta condici\'on permite reconstruir historiales, verificar intervenciones y comparar resultados productivos a lo largo del tiempo. Para MiGanado, la trazabilidad no se limita a cumplir una funci\'on documental: constituye la base necesaria para generar alertas, reportes, proyecciones y futuros modelos predictivos.

La calidad de datos tambi\'en incide en la utilidad del sistema. Registros incompletos, duplicados o inconsistentes pueden afectar indicadores, calendarios y resultados de modelos de inteligencia artificial. Por ese motivo, el dise\~no del prototipo deber\'a contemplar validaciones de carga, cat\'alogos controlados, permisos de edici\'on y mecanismos que reduzcan errores frecuentes sin dificultar el uso en contextos de campo.

\subsection{Acceso sin conexi\'on en entornos rurales}

El uso de aplicaciones digitales en establecimientos ganaderos presenta una condici\'on operativa particular: muchas tareas se realizan en zonas rurales donde la conectividad puede ser limitada, intermitente o directamente inexistente. Esta situaci\'on afecta de manera directa la posibilidad de consultar historiales, registrar eventos sanitarios y reproductivos, cargar observaciones de campo o actualizar informaci\'on de animales en el momento en que ocurren las actividades.

En este contexto, el funcionamiento sin conexi\'on constituye un requisito relevante para la experiencia de productores, veterinarios y empleados rurales. Una herramienta que dependa permanentemente de internet puede obligar a diferir la carga de datos hasta regresar a una zona con se\~nal, aumentando el riesgo de omisiones, errores de transcripci\'on o p\'erdida de informaci\'on. Por el contrario, la posibilidad de registrar datos localmente durante el trabajo en campo permite conservar la continuidad del proceso operativo y sostener la calidad de los registros.

Para MiGanado, el soporte offline se entiende como la capacidad de permitir la carga y consulta de informaci\'on esencial aun cuando el dispositivo no disponga de conexi\'on. Una vez restablecido el acceso a internet, el sistema deber\'a sincronizar los datos pendientes con la base central, resolviendo inconsistencias y manteniendo la trazabilidad de las modificaciones. Esta caracter\'istica se vincula con la propuesta del sistema porque favorece la adopci\'on en condiciones reales de uso, especialmente en establecimientos donde las decisiones sanitarias y reproductivas se toman lejos de una oficina o de una conexi\'on estable.

\section{Estado del Arte}

\subsection{Aplicaciones de gesti\'on ganadera}

En el mercado existen diversas aplicaciones orientadas a la gesti\'on ganadera que permiten digitalizar parte de las tareas realizadas en los establecimientos rurales. Entre las soluciones relevadas se identifican plataformas como Albor Agro, FieldData, iLiveStock y VacApp, las cuales abordan distintos aspectos de la administraci\'on ganadera, tales como el registro de animales, la organizaci\'on de informaci\'on productiva, el seguimiento de tareas y, en algunos casos, la planificaci\'on mediante calendarios o funcionalidades complementarias.

A partir del an\'alisis comparativo realizado, se observa que la mayor\'ia de estas herramientas se enfoca principalmente en la gesti\'on y el registro de informaci\'on ganadera. Esto representa un avance frente a los registros manuales o dispersos, ya que permite centralizar datos vinculados con animales, lotes, actividades y eventos del establecimiento. Sin embargo, no todas las soluciones relevadas integran en una misma plataforma funcionalidades vinculadas con inteligencia artificial, an\'alisis predictivo, detecci\'on de anomal\'ias, asistentes inteligentes o administraci\'on multiestablecimiento.

El relevamiento tambi\'en muestra que algunas aplicaciones incorporan funcionalidades parciales, como calendarios integrados o uso de herramientas de inteligencia artificial, pero estas capacidades no siempre se encuentran combinadas con una propuesta integral de seguimiento sanitario, reproductivo, operativo y econ\'omico. En este sentido, el an\'alisis de competencias permite identificar oportunidades de mejora relacionadas con la integraci\'on funcional, la asistencia al usuario y el apoyo a la toma de decisiones.

La Tabla~\ref{tab:state-of-art} presenta una comparaci\'on funcional entre las aplicaciones relevadas y la propuesta de MiGanado, considerando criterios vinculados con gesti\'on ganadera, calendario integrado, uso de inteligencia artificial, an\'alisis predictivo, detecci\'on de anomal\'ias, asistente inteligente y administraci\'on multiestablecimiento.

Como se observa en el cuadro comparativo, las aplicaciones relevadas cubren principalmente funciones de gesti\'on ganadera, pero presentan limitaciones en relaci\'on con capacidades inteligentes y funcionalidades avanzadas de an\'alisis. FieldData incorpora calendario integrado y uso de inteligencia artificial para el an\'alisis de audios de WhatsApp. Sin embargo, no se identifican en las herramientas comparadas funcionalidades integradas de an\'alisis predictivo, detecci\'on de anomal\'ias, asistente inteligente y gesti\'on multiestablecimiento en conjunto.

\subsection{Diferenciaci\'on de MiGanado}

MiGanado se diferencia por proponer una plataforma orientada espec\'ificamente a la gesti\'on de establecimientos bovinos de cr\'ia en Argentina. Esta delimitaci\'on permite enfocar el dise\~no en problem\'aticas propias del sector, tales como la identificaci\'on individual mediante caravana, el seguimiento sanitario y reproductivo, la organizaci\'on de animales por lotes, la administraci\'on de m\'ultiples campos y el acompa\~namiento de distintos perfiles de usuario dentro del proceso productivo.

La propuesta busca integrar en una misma herramienta funcionalidades que suelen encontrarse separadas en distintas soluciones. Entre ellas se incluyen el registro individual de animales, la gesti\'on de lotes y establecimientos, el calendario de actividades, las alertas sanitarias y reproductivas, los historiales veterinarios, los reportes, el m\'odulo econ\'omico, el asistente inteligente y componentes de inteligencia artificial aplicados al an\'alisis de informaci\'on. Esta integraci\'on busca transformar la aplicaci\'on en una herramienta de apoyo estrat\'egico para el productor, y no \'unicamente en un sistema de registro digital.

Uno de los diferenciales principales de MiGanado es la incorporaci\'on de inteligencia artificial aplicada al contexto ganadero. El sistema contempla funcionalidades de an\'alisis predictivo, detecci\'on de anomal\'ias y asistencia mediante chatbot, con el objetivo de facilitar la interpretaci\'on de datos y acompa\~nar la toma de decisiones. Estas capacidades pueden contribuir a anticipar situaciones relevantes, detectar comportamientos fuera de lo esperado y brindar orientaci\'on al usuario a partir de la informaci\'on cargada en la plataforma.

Otro aspecto diferencial es la posibilidad de uso en campo. En establecimientos rurales, la conectividad a internet puede ser limitada o intermitente, por lo que MiGanado contempla la carga de informaci\'on sin conexi\'on y la posterior sincronizaci\'on de datos cuando el dispositivo recupere conectividad. Esta caracter\'istica resulta clave para que la herramienta pueda utilizarse durante recorridas, controles sanitarios, tareas reproductivas o registros realizados directamente en el campo.

Adem\'as, la plataforma incluir\'a un m\'odulo econ\'omico que utilizar\'a la informaci\'on generada por el sistema para proyectar costos, ganancias y escenarios productivos. De esta manera, MiGanado no solo permitir\'a registrar eventos pasados, sino tambi\'en analizar informaci\'on para anticipar decisiones m\'as rentables para el establecimiento.

\subsection{Conclusiones del estado del arte}

El relevamiento realizado permite identificar que las aplicaciones ganaderas existentes aportan soluciones para digitalizar la gesti\'on de animales y establecimientos. Sin embargo, el cuadro comparativo evidencia que varias de estas herramientas se concentran en funciones operativas b\'asicas, mientras que las funcionalidades avanzadas relacionadas con inteligencia artificial, an\'alisis predictivo, detecci\'on de anomal\'ias, asistentes inteligentes y gesti\'on multiestablecimiento no aparecen integradas de manera completa.

En este contexto, MiGanado se posiciona como una propuesta integral que combina gesti\'on ganadera, calendario operativo, inteligencia artificial, an\'alisis predictivo, detecci\'on de anomal\'ias, asistente inteligente, administraci\'on multiestablecimiento y m\'odulo econ\'omico. Esta combinaci\'on permite diferenciar la plataforma frente a otras soluciones relevadas y orientar el desarrollo hacia una herramienta adaptada a las necesidades de productores, veterinarios, administradores y empleados rurales.

Por lo tanto, el aporte principal de MiGanado se encuentra en integrar capacidades operativas, anal\'iticas e inteligentes dentro de una misma soluci\'on, con foco en la ganader\'ia bovina de cr\'ia en Argentina y en las condiciones reales de trabajo en establecimientos rurales.

\begin{table}[H]
	\caption{Comparaci\'on funcional de aplicaciones de gesti\'on ganadera}
	\label{tab:state-of-art}
	\centering
	\scriptsize
	\resizebox{\textwidth}{!}{%
	\begin{tabular}{|l|c|c|c|c|c|c|c|}
		\hline
		\textbf{Aplicaci\'on} & \textbf{Gesti\'on de ganado} & \textbf{Calendario integrado} & \textbf{Utiliza IA} & \textbf{An\'alisis predictivo} & \textbf{Detecci\'on de anomal\'ias} & \textbf{Asistente inteligente} & \textbf{Multiestablecimiento} \\
		\hline
		Albor Agro & S\'i & No & No & No & No & No & No \\
		\hline
		FieldData & S\'i & S\'i & S\'i & No & No & No & No \\
		\hline
		iLiveStock & S\'i & No & No & No & No & No & No \\
		\hline
		VacApp & S\'i & No & No & No & No & No & No \\
		\hline
		MiGanado & S\'i & S\'i & S\'i & S\'i & S\'i & S\'i & S\'i \\
		\hline
	\end{tabular}
	}
\end{table}

\FloatBarrier
